\chapter{问题、解决对策与反思}
\section{需求问题：从功能堆叠回到业务目标}
\subsection{问题与原因}
项目早期的一个认识偏差，是把功能数量当作完成度的主要依据。我们会关心页面是否丰富、AI 是否能够回答、主题是否足够多，却容易忽略商品不可售、支付重复提交、骑手争抢任务和角色越权等问题。这些边界不一定在一次顺利演示中出现，却决定了系统能否被解释、被验证。

初版 SRS 中，如果只写“用户可以下单”“骑手可以接单”，就把大量规则留给了实现者自行理解。不同成员会形成不同假设：有人认为取消订单只需改状态，有人认为要同时恢复库存；有人认为 AI 推荐可以直接加入购物车，有人认为应先征得用户确认。接口联调时出现的冲突，往往在需求阶段就已经埋下。
\subsection{改进与认识}
李罡老师对初版 SRS 的意见促使我们把关注点移向需求本身。我们重新梳理角色、目标、业务边界、正常与异常流程，并为重要规则增加验收描述。对于骑手配送，补充完整履约过程；对于 AI，明确它只能辅助查找和表达，不能取得超出用户本身的权限；对于优惠，明确计算顺序和适用边界。

这一过程让我们理解，SRS 的价值不是把产品介绍写得更长，而是减少团队对同一件事的不同理解。需求越明确，实现时需要临时猜测的内容越少，测试也越容易判断对错。课程设计因此成为对问题分析、方案选择和结果验证的综合训练。
\section{按端开发还是全栈开发：对答辩问题的回应}
\subsection{两种方式的价值与局限}
答辩时，老师们询问我们认为按端开发更好，还是采用全栈开发更好。这个问题提醒我们，分工方式会直接影响成员对业务的理解、接口协同和最终交付质量，不能只按四个人对应四个页面入口进行分配。

这里的“按端开发”主要指顾客端、商家端、骑手端和管理员端分别由不同成员负责。它的优势是页面责任清楚、角色体验相对完整、成员容易形成对某类用户的理解。但订单、库存、支付和通知是跨端共享的，一旦各端分别实现自己的规则，就可能产生重复逻辑和状态冲突。

“全栈开发”在本项目中更适合理解为按业务用例纵向交付：同一成员围绕某个需求，贯通页面、接口、数据与测试。它能缩短问题定位链路，让开发者理解从输入到持久化结果的全过程；局限是成员需要更广的技术能力，也更容易同时修改公共文件。如果缺少共同规范，同样会增加冲突。
\begin{longtable}{L{.16\linewidth}L{.35\linewidth}L{.40\linewidth}}
\caption{两种分工方式的项目适用性分析}\\\toprule 维度&按角色端主责&按业务用例全栈交付\\\midrule\endfirsthead\toprule 维度&按角色端主责&按业务用例全栈交付\\\midrule\endhead
责任清晰度&页面归属直接，但跨端流程容易缺少总负责人&每个需求从界面到验证均有明确责任。\\
业务理解&容易深入单一角色，也可能局限在本端&能够理解完整数据流和失败影响。\\
协作成本&跨端接口较多，状态变更需要频繁沟通&公共服务与数据库变更需要严格协调。\\
能力要求&可结合成员擅长的页面和技术方向&需要掌握前后端基础并获得同伴评审。\\
主要风险&各端独立完成却无法组成交易闭环&重复修改基础设施，出现风格与规则差异。\\
适用方式&保持角色体验与页面主责&适用于有明确验收标准的业务纵向任务。\\\bottomrule
\end{longtable}
\subsection{本项目更合适的协作方式}
对于四人、三周的课程实践，我们认为更合理的方式是“角色端主责与业务用例贯通相结合”。每个端有人维护整体体验；每个跨端需求另有负责人，从规则、接口到测试跟踪结果；计价、权限、状态机、响应规范和数据库变更则由全组共同评审。

例如，“外送订单完成”不能拆成顾客写一个按钮、商家写一个按钮、骑手再写一个按钮后就视为完成。应先约定订单和任务状态、接口输入输出与异常，再分别实现页面与服务，最后用同一订单跑完四端流程。负责该用例的成员需要解释整条链路，但不意味着所有代码都由一个人完成。

这种分工也有利于课程学习。每名成员应当能讲清自己负责模块的需求、接口、数据和测试；公共业务通过结对评审让其他成员至少理解关键约束，避免一个成员离开后就无人能够解释核心逻辑。
\subsection{任务与协作记录}
本组采用基础技术方向与扩展模块相结合的分工方式。四名成员分别承担前端、后端、数据库和测试的主要职责，同时负责看板、红包积分、AI 与高德地图等扩展内容。表\ref{tab:team}依据小组确认的实际分工整理，不以模块数量推定工作量比例。
\begin{longtable}{L{.15\linewidth}L{.16\linewidth}L{.16\linewidth}L{.42\linewidth}}
\caption{小组实际分工}\label{tab:team}\\\toprule 成员&基础方向&扩展模块&协作接口与联调关注点\\\midrule\endfirsthead\toprule 成员&基础方向&扩展模块&协作接口与联调关注点\\\midrule\endhead
范文丽&前端&看板&对齐后端响应与页面状态，明确看板统计口径和数据展示。\\
李承文&后端&红包积分&对齐订单、优惠券与积分规则，关注计价、支付和取消回退。\\
赵紫怡&数据库&AI&协调业务数据与查询，关注真实菜单、订单授权和 AI 调用边界。\\
仲禹潼&测试&高德地图&组织功能与回归验证，关注定位、导航及地图失败后的替代路径。\\\bottomrule
\end{longtable}
这一安排结合了成员的技术主责与具体功能任务，但也意味着一个需求往往需要多人共同完成。以 AI 点餐为例，候选商品依赖数据库与 AI 模块，展示和确认依赖前端，加入购物车仍需后端执行规则，最后由测试检查完整流程。高德地图同样需要页面入口、后端数据与异常处理共同配合。

因此，测试由专人主责不意味着其他成员可以不验证自己的代码，前端主责也不意味着只需要了解页面。结合答辩意见，我们需要在保留技术主责的同时，加强对完整业务链路的理解，让每项需求都有明确的联调与验收责任。

协作记录应围绕可交付结果组织。一个任务至少说明需求依据、接口契约、涉及的数据、验收步骤和失败处理；一次提交应让其他成员看懂修改目的；一次合并应确认相关测试和跨端影响。文件数量或提交数量只能提供活动线索，不能直接等同于贡献质量。

多人协作中的沟通也不只是同步“我写完了”。更有价值的信息包括新增了什么约束、哪些接口行为改变、哪些旧假设失效、谁需要同步页面或测试。只有这些信息能够到达相关成员，分工才真正提高整体效率。
\section{AI 一次回答调用三次：对答辩问题的回应}
\subsection{先明确三次调用各自解决什么问题}
老师们在答辩中提出，AI 回答一个问题需要调用三次是否过多，在并发量增大时会出现什么问题，应如何解决。这一问题使我们认识到，能得到回答只是功能成立的起点，还需要说明调用链路为什么这样设计，以及每次调用是否提供不可替代的价值。

如果一个方案将意图识别、推荐生成和语言组织分别交给模型，三次调用可能有其设计原因，但不能默认所有请求都需要经过完整链路。平台规则咨询可以通过本地规则回答，订单状态由数据库提供，明确预算下的候选筛选也可以先在本地完成。对于简单问题，额外调用可能主要增加成本与不确定性。

因此，我们的判断不是仅凭“三次”就断言设计错误，而是逐步检查每次调用的输入、输出、依赖、失败处理和可替代方案。当前实现路径已在第四章说明；本节仍需通过请求追踪确认不同场景的真实调用次数，避免把设计图中的步骤数等同于实际模型请求数。
\subsection{并发放大、延迟与失败传播}
设用户请求到达率为 $\lambda$，每个用户请求平均触发 $k$ 次首次模型调用，则不计重试时，外部模型请求率约为：
\[
\lambda_{\mathrm{model}}=\lambda\,\mathbb{E}[k].
\]
若每个问题固定串行调用三次，各次耗时分别为 $t_1,t_2,t_3$，本地处理耗时为 $t_{\mathrm{local}}$，则无排队条件下端到端耗时约为：
\[
T=t_1+t_2+t_3+t_{\mathrm{local}}.
\]
在稳定运行、尚未过载的假设下，可以用平均到达率乘平均等待时间估算在途用户请求量。\textbf{下面仅为说明风险的假设算例，不是项目实测：}若每秒到达 20 个用户问题，每题调用 3 次、每次平均耗时 2 秒，则外部调用约为 60 次/秒；不计本地开销时，每题约需 6 秒，在途用户请求约为 120 个。若每题只需 1 次同等时长调用，则相应约为 20 次/秒、2 秒和 40 个在途用户请求。

并发增大后，问题不只体现为等待变长。外部服务可能触发请求数或令牌限额，应用线程和连接池持续占用，排队增加尾部延迟，超时后重试又进一步放大外部请求。若聊天与下单共享有限资源，AI 的拥塞还可能影响核心订单服务。

若完整答复要求三次调用全部成功，并且暂时假设各次成功概率相同且相互独立，单次成功概率为 $p$，则整条链路成功概率为 $p^3$。例如 $p=0.99$ 时约为 $97.03\%$。真实调用可能存在相关性和降级路径，因此该计算仅用于说明依赖步骤增加会扩大失败面，不能直接预测实际成功率。
\subsection{改进顺序与具体措施}
\begin{enumerate}
\item \textbf{先减少不必要的调用。}采用本地意图分流，规则问题直接答复；订单问题先校验归属并查询数据库；推荐先按预算、库存和偏好筛选。需要自然语言说明时，再对已获得的真实数据调用模型。
\item \textbf{合并能够共享上下文的步骤。}在质量可接受的前提下，让一次受约束请求返回文本与结构化候选，避免为了重新组织同一批信息再次调用。输出仍需服务端校验，不能把“减少调用”变成放宽交易规则。
\item \textbf{限制并发与排队。}设置单用户频率限制、全局 AI 并发上限和有界等待队列。达到上限时及时返回繁忙提示或本地结果，不能无限积压。AI 外部请求使用独立连接池或执行资源，保护下单和支付。
\item \textbf{设置整个回答的时间预算。}总超时覆盖首次调用、退避和重试。仅对适合重试的暂时性失败执行有限重试，并考虑服务端限流提示与抖动退避；不能在多个层级分别重试，造成乘法放大。
\item \textbf{缓存稳定内容并避免越权复用。}平台常见规则可以短期缓存；个性化内容需要按权限、偏好和数据版本区分。订单状态与库存不能盲目复用过期答案，加入购物车和支付前必须重新检查。
\item \textbf{失败时降级并保持用户知情。}模型连续失败时暂时停止无效调用，返回本地规则、候选商品或普通搜索入口。流式显示可以改善等待体验，但不会自动降低调用次数或令牌费用，不能替代容量治理。
\item \textbf{以观测结果验证优化。}记录每题调用数、重试数、耗时分位数、令牌消耗、429 比例、降级率和候选有效率，结合用户确认率评估，避免只追求请求数量减少而牺牲结果质量。
\end{enumerate}
\subsection{已实现基础与后续工作}
现有实现基础见第四章。就本节问题而言，当前客户端同步等待外部结果，仍需关注请求线程占用；全局并发预算、跨实例限流、缓存失效和完整监控未获得本次核对的充分运行证据，应列入后续改进与验收。

特别需要区分重试参数与业务编排调用数。默认最多 3 次重试不代表一题固定请求三次，更不代表所有失败都值得重试。我们应通过场景日志给出具体证据，再讨论是否合并调用或收紧预算。
\section{典型工程问题与处理思路}
\subsection{历史订单随商品修改而变化}
问题根源是混淆了当前商品数据与交易历史。当前实现通过订单明细和地址快照进行隔离。后续回归应覆盖名称、价格、地址修改与商品下架等情况，确保历史页面不重新从可变资源中拼接错误事实。由此得到的认识是：数据建模必须先确定字段表示哪个时间点的事实。
\subsection{前端金额与服务端结果不一致}
前端可以进行即时试算，但不应成为计价权威。当前报价服务集中处理真实价格、满减、会员与配送费，结算服务处理优惠券和积分。后续需要统一 SRS、接口字段和页面的展示顺序，减少用户看到的金额与最终应付不一致。金额计算应使用明确精度与舍入规则，不能简单照搬参考报告中的字符串加法方案。
\subsection{重复请求与并发抢单}
重复点击支付、网络超时重发以及多个骑手同时接单，都会让“先查询再修改”的简单逻辑暴露竞争窗口。当前实现使用幂等键、条件更新和事务作为防护基础。测试必须进一步检查受影响行数、状态记录以及资产副作用次数，不能只观察最后显示哪位骑手。
\subsection{接口契约与错误反馈不一致}
某个接口返回成功，未必意味着前端能正确解释数据。响应层级、字段名称、状态码和错误消息不一致，都可能导致页面出现空数据或错误提示。应通过共享接口定义和契约测试减少差异，并把业务失败保留为失败，不能由前端统一包装成成功。
\subsection{通知与业务事实不同步}
通知用于提醒，订单与任务记录才是权威事实。若推送断连，页面应能通过重新查询恢复状态；若通知保存失败，也需要明确错误传播与补偿方式。对这类问题，我们认识到“消息发出”与“业务完成”是不同事件，不能互相代替。
\Needspace{8\baselineskip}
\section{SRS 与当前实现的差异及改进清单}
\begin{longtable}{L{.19\linewidth}L{.36\linewidth}L{.36\linewidth}}
\caption{需要继续收敛的需求与实现差异}\\\toprule 事项&当前核对情况&改进方向\\\midrule\endfirsthead\toprule 事项&当前核对情况&改进方向\\\midrule\endhead
订单阶段命名&SRS 使用“制作中”；枚举使用“待派单”等状态&形成一份业务名、编码、页面文案的映射表，并同步测试。\\
配送到店阶段&SRS 将到店视为时间记录；当前实现有独立到店状态&明确以哪一版本为准，更新状态图与异常恢复规则。\\
计价顺序&源码含商家满减，且会员计算基于满减后金额&SRS 补充完整顺序、适用金额与舍入示例。\\
优惠券处理时机&SRS 提及下单预占；已读提交路径未见同等优惠券预占，结算时校验核销&明确锁定或支付时竞争的产品选择，补并发用券测试。\\
接口版本&新旧接口前缀并存&列出兼容范围与迁移计划，避免重复写入口绕过统一规则。\\
响应可追踪性&SRS 要求请求标识；统一响应对象本身未见该字段&核对是否由响应头提供，否则补齐统一追踪方案。\\
AI 请求预算&已有总超时和回退；分布式并发治理未获验证&增加调用预算、限流、隔离和压测证据。\\
过程与验收证据&当前快照可检查测试，完整 TDD 时序与交叉验收尚需归档&保存真实提交、运行日志和复测记录，不使用模拟材料替代。\\\bottomrule
\end{longtable}
这些差异说明项目仍有改进空间，也提供了进一步收敛的具体方向。我们希望后续每一个“已完成”的结论，都能够由需求、实现和验证共同支撑。
\section{AI 辅助开发的核查问题与改进}
AI 的两类使用方式和人机职责见第七章。这里仅记录已经暴露的核查问题：开发记录不完整时，不能从源码风格推断使用过哪些工具、采纳了哪些输出；生成的代码或文档即使表面完整，也可能把未经确认的推断写成事实，或让同一错误假设同时进入实现与测试。

后续应保留必要的开发记录和材料来源，由成员逐项审阅需求、事实、分工和答辩表述。涉及订单状态、金额、权限、数据库事务和测试预期时，成员必须核对规则与实现并承担最终判断责任，不能以 AI 输出或测试通过代替业务依据。
\clearpage
\section{课程设计的收获与不足}
通过本次实践，我们对课程设计的理解从“完成一个外卖网站”逐步转向“完成一次有依据、有协作、有验证的软件开发过程”。页面让需求可见，接口使协作明确，数据库保留事实，测试约束行为，而文档让不同成员能够围绕同一目标讨论。

我们仍然存在需求与实现同步不够及时、部分边界验证不充分、分工记录不够清晰以及对 AI 并发成本考虑不够深入等不足。对这些问题的认识，不应停留在笼统的“经验不足”，而应落实为更新状态映射、完善计价示例、补齐并发与权限测试、核对真实调用次数和保留过程证据。

这也是本次课程设计最重要的学习结果：能够运行的系统是起点，能够解释其需求、理解其限制、验证其行为，并在反馈中持续改进，才体现了软件工程实践的完整性。
